% ==============================================================================
% 10_retex_amelioration.tex — UE SIS589
% ==============================================================================
\chapter{Retour d'Expérience et Amélioration Continue}
\label{chap:retex}

\epigraph{\itshape « Un incident mal analysé est un incident
  qui se répétera. »}
         {--- Principe du RETEX opérationnel}

\begin{boxobjectifs}
\begin{enumerate}
  \item Structurer et conduire un RETEX post-incident en distinguant
        analyse des causes et recommandations.
  \item Appliquer le cycle d'amélioration continue (PDCA) à la sécurité
        opérationnelle.
  \item Définir et calculer les métriques SOC de performance à long terme.
  \item Comprendre le principe et la valeur des exercices Purple Team.
  \item Rédiger un plan de remédiation priorisé exploitable par le management.
\end{enumerate}
\end{boxobjectifs}

% ==============================================================================
\section{Le RETEX~: structure et méthode}
\label{sec:retex}
% ==============================================================================

\begin{boxdef}
\textbf{RETEX (Retour d'Expérience)~:} Processus structuré d'analyse
postérieure à un incident, visant à comprendre les causes profondes
(non seulement les symptômes), à documenter les leçons apprises et à
produire des recommandations actionnables pour éviter la récidive et
améliorer le dispositif de détection et de réponse.
\end{boxdef}

Le RETEX n'est pas un tribunal. Son objectif est la compréhension et
l'amélioration, pas la désignation d'un coupable. Un RETEX conduit dans
un climat de culpabilisation sera incomplet (les acteurs minimisent leurs
erreurs) et contre-productif.

\subsection{Structure d'un RETEX}

Un RETEX complet comporte six sections~:

\begin{enumerate}
  \item \textbf{Contexte et chronologie.}
        Résumé de l'incident~: type, dates, systèmes affectés, durée.
        Timeline consolidée (cf. chapitre~2).
  \item \textbf{Analyse des causes profondes (\textit{Root Cause Analysis}).}
        Identification des causes techniques, organisationnelles et
        humaines. Utilisation de la méthode des \og{}5~Pourquoi\fg{} ou
        du diagramme d'Ishikawa.
  \item \textbf{Analyse de la détection.}
        Pourquoi l'incident n'a-t-il pas été détecté plus tôt~?
        Quelles règles manquaient~? Quels journaux n'étaient pas collectés~?
        Le FPR était-il trop élevé~?
  \item \textbf{Analyse de la réponse.}
        La réponse a-t-elle été conforme aux playbooks~?
        Les délais SLA ont-ils été respectés~?
        Les outils étaient-ils disponibles et opérationnels~?
  \item \textbf{Recommandations.}
        Actions correctives priorisées (court, moyen, long terme)
        avec responsable, délai et indicateur de vérification.
  \item \textbf{Validation et capitalisation.}
        Validation des recommandations par le management~;
        intégration dans la documentation (playbooks, règles SIEM,
        politique de sécurité).
\end{enumerate}

\subsection{Analyse des causes profondes}

La méthode des \textbf{5~Pourquoi} est la plus accessible pour les
incidents de sécurité~:

\begin{table}[H]
\centering\small
\caption{Méthode 5~Pourquoi appliquée à l'incident LiZbiZ-SIS}
\label{tab:5pourquoi}
\rowcolors{2}{TableauPair}{TableauImpair}
\begin{tabular}{p{1cm}p{5.5cm}p{7cm}}
\toprule
\tableauHeader{\#} & \tableauHeader{Question} & \tableauHeader{Réponse} \\
\midrule
1 & Pourquoi l'attaquant a-t-il réussi à accéder au serveur~? &
Le compte \texttt{backup} avait un mot de passe faible \\
2 & Pourquoi ce compte avait-il un mot de passe faible~? &
La politique de mot de passe ne couvrait pas les comptes de service \\
3 & Pourquoi la politique ne les couvrait-elle pas~? &
La PSSI n'avait pas été mise à jour depuis 3~ans \\
4 & Pourquoi la PSSI n'a-t-elle pas été mise à jour~? &
Aucun processus de revue annuelle n'était formalisé \\
5 & Pourquoi n'y avait-il pas de processus de revue~? &
Absence de rôle RSSI à temps plein dans l'organisation \\
\midrule
\multicolumn{2}{l}{\textbf{Cause racine~:}} &
Absence de gouvernance sécurité formelle (pas de RSSI, PSSI obsolète) \\
\bottomrule
\end{tabular}
\end{table}

% ==============================================================================
\section{Le cycle PDCA appliqué à la sécurité opérationnelle}
\label{sec:pdca}
% ==============================================================================

\begin{boxdef}
\textbf{PDCA (\textit{Plan-Do-Check-Act})~:} Cycle d'amélioration continue
(Deming) appliqué à la sécurité opérationnelle. Chaque incident,
chaque exercice, chaque revue de métriques est une occasion de
progresser dans le cycle.
\end{boxdef}

Application au contexte SOC~:

\begin{description}
  \item[\textbf{Plan (Planifier)~:}] Définir les objectifs de sécurité~;
        identifier les risques prioritaires~; planifier les améliorations
        issues du RETEX~; mettre à jour les playbooks.
  \item[\textbf{Do (Faire)~:}] Déployer les améliorations~:
        nouvelles règles SIEM, formations, patches, durcissement.
  \item[\textbf{Check (Vérifier)~:}] Mesurer l'efficacité des améliorations~:
        métriques SOC (MTTD, MTTR, FPR), tests de pénétration, Purple Team.
  \item[\textbf{Act (Agir)~:}] Corriger ce qui ne fonctionne pas~;
        capitaliser ce qui fonctionne~; démarrer un nouveau cycle.
\end{description}

% ==============================================================================
\section{Métriques de performance à long terme}
\label{sec:metriques-lt}
% ==============================================================================

Au-delà des métriques opérationnelles immédiates (MTTD, MTTR,
chapitres~1 et~8), l'amélioration continue requiert des indicateurs
de progression à moyen terme~:

\begin{table}[H]
\centering\small
\caption{Métriques de performance SOC à long terme}
\label{tab:metriques-lt}
\rowcolors{2}{TableauPair}{TableauImpair}
\begin{tabular}{p{4.5cm}p{4cm}p{5cm}}
\toprule
\tableauHeader{Métrique} & \tableauHeader{Formule / Mesure} &
\tableauHeader{Objectif et interprétation} \\
\midrule
Taux de détection &
$\frac{\text{incidents détectés}}{\text{incidents totaux}} \times 100$ &
Mesure la couverture réelle~; augmenter vers 100\% \\
Couverture ATT\&CK &
$\frac{\text{techniques couvertes}}{\text{techniques totales}} \times 100$ &
Mesurer par tactique~; prioriser les angles morts \\
Réduction du MTTD &
$\text{MTTD}_n - \text{MTTD}_{n-1}$ &
Tendance décroissante = amélioration \\
Réduction du FPR &
$\frac{\text{FP}}{\text{alertes totales}}$ par période &
Tendance décroissante = règles mieux calibrées \\
Taux de récidive &
Incidents du même type sur 12~mois &
0 = recommandations RETEX bien appliquées \\
Couverture de logging &
$\frac{\text{sources actives}}{\text{sources cibles}}$ &
Mesure les angles morts de collecte \\
\bottomrule
\end{tabular}
\end{table}

Ces métriques doivent être calculées sur des périodes comparables
(mensuel, trimestriel) et présentées en tendance, pas en valeur
absolue. Un tableau de bord SOC complet est mis à jour mensuellement
et présenté au RSSI et à la direction.

% ==============================================================================
\section{Le Purple Team}
\label{sec:purple-team}
% ==============================================================================

\begin{boxdef}
\textbf{Red Team~:} Équipe offensive qui simule un attaquant réel pour
tester les défenses de l'organisation. Cherche à atteindre des objectifs
définis (accès à la base de données clients, élévation de privilèges
sur l'AD) en utilisant les mêmes TTPs qu'un vrai attaquant.
\end{boxdef}

\begin{boxdef}
\textbf{Blue Team~:} Équipe défensive~: le SOC. Détecte, analyse et
répond aux actions du Red Team (et aux vrais incidents).
\end{boxdef}

\begin{boxdef}
\textbf{Purple Team~:} Mode collaboratif Red Team + Blue Team, opérant
\textbf{ensemble} pour maximiser la valeur de l'exercice. Le Red Team
exécute une technique~; le Blue Team vérifie s'il l'a détectée~; si non,
ils analysent ensemble pourquoi et corrigent la règle ou la source de
logs manquante. L'objectif est l'amélioration de la détection,
pas la démonstration de compétences offensives.
\end{boxdef}

\subsection{Déroulement d'un exercice Purple Team}

\begin{enumerate}
  \item \textbf{Planification.} Sélectionner les TTPs à tester
        (priorité~: angles morts de la couverture ATT\&CK).
        Définir le périmètre, les systèmes cibles, la fenêtre temporelle.
  \item \textbf{Exécution.} Le Red Team exécute la technique choisie.
        Le Blue Team surveille en temps réel.
  \item \textbf{Vérification.} Les deux équipes confrontent leurs
        observations~: l'alerte a-t-elle été générée~?
        Dans quel délai~? Était-elle actionnable~?
  \item \textbf{Remédiation.} Si la technique n'a pas été détectée~:
        identifier la source de logs manquante ou la règle à créer.
        Créer la règle. Re-tester immédiatement.
  \item \textbf{Documentation.} Chaque TTP testé est documenté dans
        le registre de couverture ATT\&CK avec le résultat~:
        Détecté / Non détecté / Partiellement détecté.
\end{enumerate}

\begin{lstlisting}[style=StyleTerminal,
  caption={Exemple de scénario Purple Team --- T1003.001 LSASS Dumping},
  label={lst:purple-team}]
# ─── RED TEAM : Dumper les credentials depuis LSASS ─────────────────────────
# (sur une machine de test en lab, avec autorisation explicite)
# Méthode 1 : Task Manager (GUI) — clic droit sur lsass.exe > Create Dump File
# Méthode 2 : procdump (Sysinternals)
procdump.exe -ma lsass.exe C:\Temp\lsass.dmp

# Méthode 3 : comsvcs.dll (LOLBin, T1218)
rundll32.exe C:\Windows\System32\comsvcs.dll MiniDump \
  (Get-Process lsass).id C:\Temp\lsass.dmp full

# ─── BLUE TEAM : Vérification de la détection ────────────────────────────────
# Question 1 : Wazuh / Sysmon a-t-il généré une alerte ?
# Chercher dans Kibana :
event.code: "10"       # Sysmon EID 10 : ProcessAccess sur LSASS
  AND winlog.event_data.TargetImage: "*lsass*"

# Question 2 : L'accès à LSASS est-il loggué en Event ID 4656 ?
event.code: "4656" AND winlog.event_data.ObjectName: "*lsass*"

# ─── SI NON DÉTECTÉ : remédiation immédiate ──────────────────────────────────
# Activer Sysmon EID 10 dans sysmonconfig.xml
# <ProcessAccess onmatch="include">
#   <TargetImage condition="end with">lsass.exe</TargetImage>
# </ProcessAccess>

# Ajouter règle Wazuh pour EID 10 + LSASS
# Re-tester : la même technique doit maintenant générer une alerte

# ─── DOCUMENTER le résultat ──────────────────────────────────────────────────
echo "T1003.001 | 2026-04-20 | Non détecté | Sysmon EID10 absent | Corrigé 2026-04-20 | Re-testé OK" \
  >> /var/log/purple_team_registry.csv
\end{lstlisting}

% ==============================================================================
\section{Plan de remédiation priorisé}
\label{sec:plan-remediation}
% ==============================================================================

Le RETEX produit des recommandations qui doivent être transformées
en plan d'actions opérationnel. Le plan de remédiation répond aux
questions~: quoi faire, dans quel ordre, par qui, pour quand,
comment vérifier.

\subsection{Priorisation}

La priorisation des actions repose sur deux axes~:
\begin{itemize}
  \item \textbf{Impact~:} réduction du risque résiduel si l'action est
        implémentée (critique, haute, moyenne, basse).
  \item \textbf{Effort~:} coût d'implémentation (humain, financier,
        délai). Une action à fort impact et faible effort est
        une \og{}quick win\fg{} à traiter en priorité.
\end{itemize}

\begin{table}[H]
\centering\small
\caption{Plan de remédiation post-incident LiZbiZ-SIS (extrait)}
\label{tab:plan-remediation}
\rowcolors{2}{TableauPair}{TableauImpair}
\begin{tabular}{p{4cm}p{2cm}p{2cm}p{2cm}p{3.5cm}}
\toprule
\tableauHeader{Action} & \tableauHeader{Impact} &
\tableauHeader{Effort} & \tableauHeader{Délai} &
\tableauHeader{Indicateur de vérification} \\
\midrule
Activer MFA sur tous les accès SSH & \gravite{critique} & Faible &
J+7 & Rapport d'audit SSH \\
Segmenter le réseau backup & \gravite{haute} & Moyen &
J+30 & Test de pénétration \\
Désactiver les comptes de service inutilisés & \gravite{haute} & Faible &
J+15 & Revue AD mensuelle \\
Déployer Sysmon sur tous les serveurs & \gravite{haute} & Faible &
J+14 & Couverture agent Wazuh 100\% \\
Mettre à jour la PSSI + politique MDP & \gravite{moyenne} & Élevé &
J+90 & Validation direction + diffusion \\
Former les équipes IT à l'IR & \gravite{moyenne} & Moyen &
J+60 & Attestation de formation \\
\bottomrule
\end{tabular}
\end{table}

\subsection{Suivi et gouvernance}

Le plan de remédiation doit être suivi~:
\begin{itemize}
  \item Revue mensuelle des actions en cours avec le RSSI.
  \item Rapport trimestriel à la direction (état d'avancement,
        métriques SOC, incidents survenus).
  \item Revue annuelle de la PSSI et des playbooks.
  \item Test annuel du plan de continuité d'activité.
\end{itemize}

\begin{boxlizbiz}
\textbf{LiZbiZ-SIS~: RETEX de l'incident du 12/04/2026}

\textbf{Cause racine~:} Absence de gouvernance sécurité (PSSI obsolète,
pas de RSSI, politique de mots de passe incomplète).

\textbf{Facteurs contributifs~:} MTTD $>$ 21~jours (aucun SIEM)~;
compte de service avec mot de passe faible non couvert par la politique~;
absence de monitoring SSH nocturne~; pas de segmentation du réseau backup.

\textbf{Bonnes pratiques observées~:} Les journaux du firewall pfSense
étaient conservés 90~jours et ont permis la reconstruction de la
timeline~; les backups n'ont pas été chiffrés (pas de ransomware).

\textbf{5~recommandations prioritaires~:} (1)~MFA sur tous les accès
distants~; (2)~Déploiement Wazuh + ELK immédiat~; (3)~Désactivation
des comptes de service inactifs~; (4)~Politique de mot de passe
couvrant tous les comptes (humains ET services)~;
(5)~Nomination d'un RSSI (même à temps partiel).
\end{boxlizbiz}

% ==============================================================================
\section*{Résumé}
% ==============================================================================

\begin{boxresume}
\begin{itemize}
  \item RETEX~: 6~sections (contexte, causes profondes, détection,
        réponse, recommandations, capitalisation). Méthode 5~Pourquoi
        pour la cause racine.
  \item PDCA appliqué au SOC~: Plan~$\to$~Do~$\to$~Check~$\to$~Act.
        Chaque incident alimente le cycle.
  \item Métriques à long terme~: taux de détection, couverture ATT\&CK,
        évolution MTTD/MTTR, taux de récidive.
  \item Purple Team~: collaboration Red + Blue pour améliorer la détection
        TTP par TTP. Documenter chaque test dans le registre de couverture.
  \item Plan de remédiation~: prioriser sur impact $\times$ effort.
        Quick wins en premier, gouvernance en parallèle.
\end{itemize}
\end{boxresume}

\begin{boxexo}
\textbf{Ex.~10.1 — RETEX structuré.}
Rédigez le RETEX complet de l'incident LiZbiZ-SIS du 12/04/2026
en suivant les 6~sections du modèle. La section causes profondes
doit utiliser les 5~Pourquoi. La section recommandations doit
contenir au moins 6~actions priorisées.

\medskip\textbf{Ex.~10.2 — Tableau de couverture ATT\&CK.}
Pour les 7~techniques MITRE identifiées dans l'incident
(tableau~\ref{tab:mitre-mapping}), remplissez un tableau de couverture~:
technique, source de log requise, règle SIEM existante (oui/non),
action si non existante.

\medskip\textbf{Ex.~10.3 — Scénario Purple Team.}
Planifiez un exercice Purple Team pour tester 3~techniques choisies
dans les angles morts de la couverture ATT\&CK de LiZbiZ-SIS~:
description de l'attaque simulée, source de log attendue, règle SIEM
à créer ou vérifier.
\end{boxexo}
